\section{参数、搜索范围与复现条件}\label{app:parameters}
原基准的八种表示由四个列集合乘有无PCA组成；八项相关列为AST、总胆固醇、红细胞压积、血红蛋白、血小板比积、红细胞平均血红蛋白量、白球比例及中性粒细胞百分比。表示比较先固定基准参数，参数搜索后续进行，不声称在表示与所有参数的完整笛卡尔积上穷举。

原基准候选范围见表\ref{tab:baseline-search}。Ridge惩罚100在本轮预先固定。网络的两个表示及轮数由训练内部验证选择，不列入下表的传统模型参数网格。
\begin{table}[htbp]
\small
\setlength{\tabcolsep}{5pt}
\centering
\caption{原基准参数候选范围与抽取数量}
\label{tab:baseline-search}
\begin{tabular}{L{0.17\linewidth}L{0.64\linewidth}L{0.09\linewidth}}
\toprule
模型 & 实际候选空间 & 抽取数 \\
\midrule
决策树 & 深度3/4/5/6/7；叶最少10/20/40/60；节点分裂最少2/10/20 & 24 \\
AdaBoost & 轮数100/200/300；学习率0.005/0.01/0.03/0.1；误差linear/square/exponential & 18 \\
原XGBoost & 轮数100/200/350；学习率0.02/0.05/0.1；深度2/3/4；子节点Hessian和10/30/60；L2=10/30/60；行和列比例各0.8/1.0 & 24 \\
随机森林 & 深度8/15/不限；叶最少7/15/30；特征比例0.6/0.8/1.0；树数200固定 & 12 \\
\bottomrule
\end{tabular}
\end{table}


联合搜索范围见表\ref{tab:joint-search}。对数均匀分布表示在对数尺度等概率抽样；整数区间注明实际包含端点。各配置使用预定随机种子抽样，不根据B误差改变抽样方案。
\begin{table}[htbp]
\small
\setlength{\tabcolsep}{5pt}
\centering
\caption{联合搜索的预声明参数范围}
\label{tab:joint-search}
\begin{tabular}{L{0.15\linewidth}L{0.66\linewidth}L{0.09\linewidth}}
\toprule
家族 & 范围与固定处理 & 抽取数 \\
\midrule
XGBoost & 轮数整数200--1000；深度整数2--6；学习率对数均匀[0.01,0.12]；子节点Hessian和对数均匀[5,120]；L2对数均匀[0.1,150]；L1=0/0.1/1/10；gamma=0/0.05/0.2/1；行、列比例各0.65/0.8/1；填补/原生缺失二选一 & 100 \\
LightGBM & 轮数250/500/800；学习率0.015/0.03/0.06；叶数7/15/31；叶最少20/40/80；L2=1/10/50；列比例0.7/1；原生缺失 & 12 \\
HistGB & 轮数200/400/700；学习率0.02/0.05；叶数7/15；叶最少20/50/80；L2=1/10/50；关闭早停、原生缺失 & 8 \\
\bottomrule
\end{tabular}
\end{table}


原大网络首个Dense后使用Dropout0.4，其余三个隐藏Dense分别接BatchNormalization与Dropout0.6/0.5/0.6；后四个Dense使用原构造中的normal初始化。小网络采用ReLU隐藏层、隐藏核L2=0.001和逐层Dropout，参数由式\eqref{eq:smallnn}及正文结构确定。原大网络的旧PReLU序列化需要按冻结结构重建后恢复权重，已有核验记录仅证明相应推理恢复，不代表恢复优化器状态续训。

\begingroup
\raggedright
历史运行版本为NumPy1.26.4、pandas2.0.3、SciPy1.11.1、scikit-learn1.3.0、statsmodels0.14.0、TensorFlow2.16.2、Keras3.8.0、XGBoost2.1.3、LightGBM4.1.0和factor-analyzer0.4.1。scikit-learn提供了相关回归与预处理实现\cite{sklearn}；文末1.3系列网页现标1.3.2，算法细节同时依据实际安装源码核对，不将网页版本当作实验运行版本。
\par
\endgroup
